% Generated from src/content/cv*.json by scripts/render_cv.py.
\documentclass[11pt,a4paper]{article}
\usepackage{fontspec}
\setmainfont{TeX Gyre Heros}
\setsansfont{TeX Gyre Heros}
\setmonofont{Latin Modern Mono Light}[Scale=MatchLowercase]
\renewcommand{\familydefault}{\sfdefault}
\usepackage{xeCJK}
\setCJKmainfont{Noto Sans CJK JP}
\setCJKsansfont{Noto Sans CJK JP}
\setCJKmonofont{Noto Sans Mono CJK JP}
\usepackage[top=0.6in,bottom=0.6in,left=0.7in,right=0.7in,headheight=14pt,headsep=12pt]{geometry}
\usepackage{xcolor,hyperref,needspace,fancyhdr,lastpage,fontawesome}
\definecolor{ink}{HTML}{171717}
\definecolor{muted}{HTML}{555555}
\definecolor{paper}{HTML}{F3F3F1}
\hypersetup{hidelinks,pdftitle={Ruşen Birben CV},pdfauthor={Ruşen Birben}}
\setlength{\parindent}{0pt}
\setlength{\parskip}{0pt}
\raggedright
\linespread{1.05}
\color{ink}
\pagestyle{fancy}
\fancyhf{}
\fancyhead[L]{\ttfamily{}\small{}Ruşen Birben\enspace{} / CV}
\fancyhead[R]{\ttfamily{}\small{}rusen.ai}
\fancyfoot[L]{\ttfamily{}\footnotesize{}rusen.ai/cv}
\fancyfoot[R]{\ttfamily{}\footnotesize{}\thepage\ / \pageref*{LastPage}}
\renewcommand{\headrulewidth}{0pt}
\fancypagestyle{firstpage}{\fancyhead{}}
\newcounter{cvsec}
\newcommand{\cvsection}[2]{%
  \par\needspace{7\baselineskip}\vspace{6pt}\refstepcounter{cvsec}%
  {\setlength{\fboxsep}{3pt}\colorbox{ink}{\color{white}\ttfamily{}\bfseries{}\small{}\ifnum\value{cvsec}<10 0\fi\thecvsec}}%
  \hspace{7pt}{\large{} #1}\hspace{7pt}{\ttfamily{}\bfseries{}\fontsize{12}{14}\selectfont{} #2}%
  \par\vspace{4pt}{\color{ink}\hrule height 0.9pt}\nopagebreak\vspace{7pt}%
}
\newcommand{\dates}[1]{\hfill{\ttfamily{}\small{}\color{muted}#1}}
\newcommand{\cvbody}[1]{{\small{}\raggedright #1\par}}
\newcommand{\extarrow}{\,{\scriptsize\faExternalLink}}
\begin{document}
\thispagestyle{firstpage}
\hrule height 2.5pt\vspace{6pt}
{\ttfamily{}\fontsize{30}{33}\selectfont{}\bfseries{} Ruşen Birben}\par\vspace{3pt}
{\ttfamily{}\fontsize{11}{13}\selectfont{}AI研究者・エンジニア}\par\vspace{6pt}
{\small{}
\faMapMarker\enspace{}アンカラ、トルコ\quad{}
\faEnvelope\enspace{}\href{mailto:contact@rusen.ai}{contact@rusen.ai}\quad{}
\faGlobe\enspace{}\href{https://rusen.ai}{rusen.ai}\par\vspace{4pt}
\faGithub\enspace{}\href{https://github.com/rusenbb}{github.com/rusenbb}\quad{}
\faLinkedin\enspace{}\href{https://linkedin.com/in/rusenbirben}{linkedin.com/in/rusenbirben}
}\par\vspace{6pt}
\cvbody{イスタンブール工科大学の人工知能・データ工学科を卒業し、中東工科大学の応用NLPグループで独立研究者として活動。LLMの頑健性と信頼性に関心を持ち、トルコ語NLP、機械忘却、金融分野のLLM応用に取り組んできました。}

\cvsection{\faGraduationCap}{学歴}
\begin{minipage}{\linewidth}\raggedright
\textbf{人工知能・データ工学 学士}\dates{2022 - 2025}\par
{\small{}イスタンブール工科大学}\dates{GPA: 3.61 / 4.00}\par
\end{minipage}\par\vspace{5pt}
\begin{minipage}{\linewidth}\raggedright
\textbf{電子・通信工学 学部在籍}\dates{2020 - 2022}\par
{\small{}イスタンブール工科大学}\par
{\small{}\itshape{}\color{muted}人工知能・データ工学への学内転科。}\par\end{minipage}\par\vspace{5pt}


\cvsection{\faTrophy}{受賞歴}
\begin{minipage}{\linewidth}\raggedright
\textbf{\small{}1416 YLSY 奨学金 · 修士・博士課程}\dates{2026}\par\vspace{2pt}
\cvbody{トルコ国民教育省による海外修士・博士課程向け1416 YLSY奨学金に採用。}
\end{minipage}\par\vspace{3pt}
\cvsection{\faFlask}{研究経験}
\begin{minipage}{\linewidth}\raggedright
\textbf{独立研究者}\dates{2026 - 現在}\par
{\small{}\itshape{}\href{https://nlp.ceng.metu.edu.tr}{中東工科大学 応用NLPグループ\extarrow}}\dates{アンカラ、トルコ}\par\vspace{3pt}
\cvbody{Çağrı Toramanとともに、SIGTURK 2027ではテュルク諸語と文化的知識を評価するベンチマークのデータ作成・評価手法に貢献し、LlamaTurk2の学習データ収集と事前学習実験にも参加。研究室の\href{https://www.youtube.com/@metunlpconnect}{METU NLP Connect}を運営し、専門家インタビュー、LLM・NLP解説、最新論文の紹介も担当。}
\end{minipage}\par\vspace{6pt}
\begin{minipage}{\linewidth}\raggedright
\textbf{AIOps 研究インターン}\dates{08/2024 - 09/2024}\par
{\small{}\itshape{}\href{https://havelsan.com/tr}{Havelsan\extarrow}}\dates{イスタンブール、トルコ}\par\vspace{3pt}
\cvbody{予知保全に向けたソフトウェアログ分析を研究し、HDFSログを使ったプロトタイプを開発。}
\end{minipage}\par\vspace{6pt}
\begin{minipage}{\linewidth}\raggedright
\textbf{NLP 研究インターン}\dates{07/2024 - 08/2024}\par
{\small{}\itshape{}\href{https://nlp.itu.edu.tr/}{ITU NLP 研究室\extarrow}}\dates{イスタンブール、トルコ}\par\vspace{3pt}
\cvbody{金融分野のLLM応用を研究。感情分析、ツール利用、推論、予測、検索拡張生成（RAG）を検討し、金融推論LLMを扱う卒業制作Kairaにつながる実験の一部を実施。}
\end{minipage}\par\vspace{6pt}
\cvsection{\faBriefcase}{職務経歴}
\begin{minipage}{\linewidth}\raggedright
\textbf{共同創業者・技術責任者}\dates{08/2025 - 04/2026}\par
{\small{}\itshape{}\href{https://github.com/fiction-ai-studios}{Fiction Studios\extarrow}}\dates{アンカラ、トルコ}\par\vspace{3pt}
\cvbody{ビルケント・サイバーパークで友人とFiction Studiosを共同創業し、技術責任者を担当。技術開発を主導し、創業者2人を支援する追加メンバー3人の活動を調整した。博物館・観光客向けの多言語LLMガイドを開発し、対話型の来館体験づくりと並行して顧客・投資家との協議にも参加。}
\end{minipage}\par\vspace{6pt}
\begin{minipage}{\linewidth}\raggedright
\textbf{ジュニア AI エンジニア}\dates{03/2025 - 07/2025}\par
{\small{}\itshape{}\href{https://cyberquote.com/tk/}{CyberQuote\extarrow}}\dates{イスタンブール、トルコ}\par\vspace{3pt}
\cvbody{Phillip Capital子会社CyberQuoteの新設AIチームに2人目のメンバーとして入社。ニュース感情分析パイプライン、金融専門家向け社内LLMツール、Web統合、モバイル音声アシスタントを開発。分散した情報源やPDFレポートから最新の金融情報を取得する際の難しさや、情報不足による回答の信頼性低下に取り組んだ。}
\end{minipage}\par\vspace{6pt}

\newpage
\cvsection{\faCode}{プロジェクト}
\begin{minipage}{\linewidth}\raggedright
{\ttfamily{}\bfseries{}Kaira}\dates{2025}\par
{\small{}\itshape{}\color{muted}卒業制作：金融推論LLM}\hfill{\small{}%
%
}\par\vspace{4pt}
\cvbody{イスタンブール工科大学の卒業制作として、Qwen系LLMの金融推論を研究。UnslothとLoRAを用いてGRPOでファインチューニングし、FinQAとGSM8Kで報酬設計や推論の振る舞いを調べた。Ollamaを使ったチャットアプリも開発。}
\end{minipage}\par\vspace{5pt}
\begin{minipage}{\linewidth}\raggedright
{\ttfamily{}\bfseries{}AIZheimer}\dates{2024}\par
{\small{}\itshape{}\color{muted}Stable Diffusionの機械忘却}\hfill{\small{}%
\href{https://github.com/rusenbb/ai-zheimer}{\faGithub\enspace{}コード}%
}\par\vspace{4pt}
\cvbody{Forget-Me-Notを基にStable Diffusion 2.1向けのクロスアテンションに基づく機械忘却を実装・調整。Attention controller/processorを開発し、attention re-steering、distraction、random stimulusの損失関数を検討。定性的実験と可視化を通じて、対象概念の除去とそれ以外の生成能力の維持とのトレードオフを調べた。}
\end{minipage}\par\vspace{5pt}
\begin{minipage}{\linewidth}\raggedright
{\ttfamily{}\bfseries{}Biting The Bytes}\dates{2024}\par
{\small{}\itshape{}\color{muted}トルコ語ダイアクリティカル復元}\hfill{\small{}%
\href{https://github.com/rusen-archive/Biting-The-Bytes}{\faGithub\enspace{}コード}\quad{}\href{https://huggingface.co/spaces/rusen/diacritize-tr}{\faExternalLink\enspace{}デモ}%
}\par\vspace{4pt}
\cvbody{トルコ語の文字復元をバイト単位のトークン分類として扱い、ByT5-smallをPEFT/LoRAでファインチューニング。モデル学習とハイパーパラメータ調整を主に担当し、デバッグとHugging Face Spacesへの公開はチームメイトと分担。ノイズやASCII表記を含む文章からş、ç、ı、ğ、ü、öなどの文字を復元する課題に取り組んだ。}
\end{minipage}\par\vspace{5pt}
\begin{minipage}{\linewidth}\raggedright
{\ttfamily{}\bfseries{}To-AI-or-Not-to-AI}\dates{2023}\par
{\small{}\itshape{}\color{muted}AI 生成テキスト検出器}\hfill{\small{}%
\href{https://github.com/rusen-archive/To-AI-or-Not-to-AI}{\faGithub\enspace{}コード}\quad{}\href{https://huggingface.co/spaces/rusen/gpt_detector}{\faExternalLink\enspace{}デモ}%
}\par\vspace{4pt}
\cvbody{RoBERTaベースの検出器3種をファインチューニングし、多数決で予測を統合。26,270件のデータ作成・前処理を分担し、Hugging Face Spacesへの公開を担当。同一分布から無作為に分割したテストセットで94.21\%の正解率を達成したが、他言語や大きく異なるドメインへの汎化は評価していない。}
\end{minipage}\par\vspace{5pt}
\begin{minipage}{\linewidth}\raggedright
{\ttfamily{}\bfseries{}Tobor}\dates{2024 - 2025}\par
{\small{}\itshape{}\color{muted}自律走行倉庫ロボット}\hfill{\small{}%
\href{https://github.com/rusen-archive/Amazon-Warehouse-Robot}{\faGithub\enspace{}コード}\quad{}\href{https://youtu.be/h-Jxs_y9kNA}{\faYoutubePlay\enspace{}動画}%
}\par\vspace{4pt}
\cvbody{ROS2のナビゲーションとYOLOによる認識を組み合わせた自律走行倉庫ロボットをチームで開発。箱の持ち上げ・配置機構を実装し、データのアノテーションを共同で行い、他のアルゴリズム開発も支援。}
\end{minipage}\par\vspace{5pt}

\cvsection{\faCogs}{スキル}
{\small{}\textbf{言語}\enspace{} Python / JavaScript/TypeScript / SQL / C/C++}\par\vspace{4pt}
{\small{}\textbf{フレームワーク・ツール}\enspace{} PyTorch / TensorFlow / Transformers / PEFT / LoRA / Unsloth / Scikit-learn / Pandas / NumPy / ROS2 / YOLO}\par\vspace{4pt}
{\small{}\textbf{専門分野}\enspace{} NLP / LLM / RAG / 機械忘却 / モデル評価}\par\vspace{4pt}

\cvsection{\faCertificate}{受講コース / 言語}
\textbf{\small{}\href{https://coursera.org/share/9c59e9758d378fb8a8b4ab7074dd71ed}{Deep Learning Specialization\extarrow}}\dates{DeepLearning.AI · Coursera}\par\vspace{2pt}
\cvbody{ニューラルネットワーク、CNN、系列モデル、最適化、機械学習プロジェクトの構成。}\vspace{5pt}
\textbf{\small{}\href{https://coursera.org/share/f8ab2981ab0bbdc3fec17c6e88804ff8}{AWS Cloud Technical Essentials\extarrow}}\dates{AWS · Coursera}\par\vspace{2pt}
\cvbody{クラウドアーキテクチャとEC2、Lambda、S3、RDS、DynamoDBなどのAWSサービス。}\vspace{5pt}

\vspace{5pt}
\faLanguage\enspace{}
{\small{}\textbf{トルコ語}\enspace{}母語\quad{}\textbf{英語}\enspace{}TOEFL iBT 109/120（2025年11月8日）}\par

\cvsection{\faHeartO}{趣味}
\noindent\begin{minipage}[t]{0.235\linewidth}\raggedright
{\large{}\faCoffee}\par\vspace{5pt}{\ttfamily{}\bfseries{}\small{}コーヒー}\par\vspace{4pt}
\cvbody{V60 / エアロプレス / モカポット}
\end{minipage}%
\hfill%
\noindent\begin{minipage}[t]{0.235\linewidth}\raggedright
{\large{}\faBook}\par\vspace{5pt}{\ttfamily{}\bfseries{}\small{}読書}\par\vspace{4pt}
\cvbody{哲学 / AI/CS / 金融}
\end{minipage}%
\hfill%
\noindent\begin{minipage}[t]{0.235\linewidth}\raggedright
{\large{}\faCamera}\par\vspace{5pt}{\ttfamily{}\bfseries{}\small{}\href{https://rusenbirben.com}{写真\extarrow}}\par\vspace{4pt}
\cvbody{ストリート写真 / 二重露光 / 野鳥写真}
\end{minipage}%
\hfill%
\noindent\begin{minipage}[t]{0.235\linewidth}\raggedright
{\large{}\faFilm}\par\vspace{5pt}{\ttfamily{}\bfseries{}\small{}アニメ・漫画}\par\vspace{4pt}
\cvbody{視聴 / 読書 / コレクション}
\end{minipage}%
\par
\end{document}
